% Verification methodology and toolchain (issue #20).

\section{Verification Methodology}
Verification checks both the values of $C=AB$ and the timing and handshake
contract in \path{docs/spec.md}. Tests progress from individual modules to
the array and the complete accelerator, so a failure can be traced to an
arithmetic, forwarding, skew or control boundary.

\subsection{Tool Flow}
The synthesizable Verilog-2005 RTL first passes static checks with
Verilator~\cite{verilatorguide}. The \texttt{make lint} target discovers all
RTL files and runs \texttt{--lint-only -Wall}, selecting each module as a
top and using the Verilog-2005 language setting. Warnings fail the target;
intentional unused signals have local waivers with explanations.

The \texttt{make test} target generates golden vectors, compiles each
\path{tb/tb_*.v} with Icarus Verilog using \texttt{-g2005 -Wall}, and runs
the resulting executable with \texttt{vvp}~\cite{icarussimulation}.
Testbenches dump VCD waveforms into \path{build/}. GTKWave provides manual
inspection of the operand wavefront, accumulator updates and controller
transitions~\cite{gtkwavedocs}; saved layouts are in \path{sim/waves/}.
The \texttt{make waves} target reruns the top-level test and opens its VCD
in the selected viewer. Waveform inspection supplements the automated
comparisons. Synthesis, resource counts and maximum-frequency measurement
are Phase~2 work and are not part of this simulation flow.

\subsection{Hierarchical Testbenches}
The MAC test checks signed products, including $-128\times-128$ and
$-128\times127$, an accumulation chain, reset, clear priority and holding
with enable low. The PE test adds the one-enabled-cycle forwarding delay
and checks that both forwarding registers hold and clear with the
accumulator. Input-buffer tests compare every A and B lane at each step
with the spec's skew table; output-buffer tests check parallel capture,
holding the previous result and readout at every address.

Controller tests check the state sequence, busy-state rejection of
\texttt{start}, restart from \texttt{DONE}, mid-run reset and the exact
compute and busy lengths for $N=3$ and $N=4$. Array tests drive skewed
operands at the edges and check every PE's partial sum after each compute
edge, as well as final results, stalls and restart. They exercise both
array sizes, catching row/column indexing and fill/drain errors that a
single final-result comparison could miss.

The end-to-end test loads operands through the host ports and compares
every output element with golden vectors. It also checks ignored writes
while busy, back-to-back runs, replacing only one operand matrix,
preservation of the previous $C$ until capture, and reset during a run.
For each golden case it counts busy cycles and checks
\texttt{cycle\_count}: both must equal $K+2N$, or 12 for $N=K=4$.
The CSV in \path{results/phase1/cycle_counts.csv} records these measurements;
its compute-cycle field is derived by subtracting the one CLEAR and one
CAPTURE cycle. The controller tests check the compute interval directly
against $T=K+2N-2$, which is 10 cycles at the default size.

\subsection{Golden Results and Directed Cases}
The Python reference in \path{model/golden_model.py} computes signed
matrix products with NumPy integer arithmetic and checks operand and
result ranges. A separate cycle-level model starts with zero forwarding
registers and accumulators, applies the specified input skew, and advances
operands one PE per cycle. The vector generator refuses to write its
files unless this model agrees with the mathematical product after
exactly $T$ steps. Expected results therefore come from the reference or
hand-worked values, rather than recorded RTL output.

The default vector set contains six directed cases---identity, zero,
all ones, maximum negative operands, mixed signed extremes and a known
product with $A=1,\ldots,16$ in row-major order and $B=A^{\mathsf T}$---plus
100 random cases spanning the signed 8-bit range. The golden-vector seed
is 1; array tests use seeds 1 and 2 for $N=4$ and $N=3$. Seeds are printed
in simulation logs, and the vector manifest maps cases to file offsets.
Separate Python self-tests check the reference, cycle model, hexadecimal
encoding and generator against hand-worked results and the spec.

\subsection{Failure Reporting and Reproduction}
Each self-checking testbench prints the following summary on success:
\begin{center}
\small\verb|PASS tb_<module> <passed>/<total>|
\end{center}
A successful test terminates with \verb|$finish|. A mismatch reports a
failure and calls
\verb|$fatal(1, ...)|, producing a nonzero simulator exit. Make also rejects
a \texttt{FAIL} summary or a missing, correctly formatted \texttt{PASS}
summary. Make retains compilation, lint and simulation logs in the
generated build directory. Repeated Make invocations rerun simulations even when
compiled executables already exist. Generated vectors and VCDs remain
uncommitted; selected logs, cycle counts and waveform captures in
\path{results/phase1/} provide report evidence.

\subsection{Continuous Integration and Scope}
The GitHub Actions simulation workflow runs on every pull request and on
pushes to \texttt{dev}. On Ubuntu~24.04 it selects Python~3.12, installs
Icarus Verilog, Verilator and NumPy, then executes the same
\texttt{make lint test} command used locally. Failed jobs upload available
\path{build/*.log} files as diagnostic artifacts. The Python model
self-tests, GUI viewer and report build are separate local checks; they
are not steps in this workflow.

All seven available RTL testbenches pass the current regression,
including all 106 golden cases end to end. These checks establish the
observed functional and cycle-level behavior for the exercised cases;
they are not an exhaustive proof or a measurement of physical timing.
Phase~2 synthesis will assess implementation resources and timing.
